% ===========================================================================
\section{Validation}
\label{sec:validation}
% ===========================================================================

\subsection{Against EDGAR: fiscal-year blended, category-mapped}

Matching the period comes first, because it moves the comparison more than anything else
in this section---and it cannot be done exactly. EDGAR reports calendar years; this
sector reports July 2023 to June 2024, so half of its reference period falls in each of
EDGAR's 2023 and 2024 columns. The closest annual approximation EDGAR's own data
supports is the half-and-half blend,
$0.5 \times \SI{47.50}{} + 0.5 \times \SI{48.76}{} = \SI{48.13}{\mega\tonne}$, and that
is what this paper compares against. The single-year comparisons are kept beside it,
because the spread between them \emph{is} the temporal uncertainty and should be visible
rather than chosen. Getting the period right has history here: one earlier edition
compared a fiscal year against a single calendar column and called it year-matched, and
the rc4 edition led with a \num{0.70} ratio whose engine side mixed the harmonised
FY2023--24 road figure with pre-harmonisation FY2021--22 rail and aviation rows against
a CY2021/CY2022 EDGAR blend---three vintages in one number, labelled ``all FY2021--22''.
That construction is reproduced at rank \emph{superseded} in the comparison table so it
remains checkable, and it must not be quoted as current.

Matching the scope comes second, and this revision does it from EDGAR's own source
categories rather than by assertion---though the evidence is inferred rather than
published. The GHG booklet defines Transport only as "mobile combustion (road \& rail \&
ship \& aviation)", with no IPCC split at all. EDGAR's category-resolved
\emph{air-pollutant} files do carry the Pakistani rows explicitly: road (1.A.3.b),
railways (1.A.3.c) and civil aviation (1.A.3.a) are present, there is \emph{no}
water-borne navigation row (1.A.3.d) for Pakistan at all, and other transportation
(1.A.3.e) is nil
(\texttt{data/activity/transport/edgar\_pak\_transport\_categories.csv}). Those files are
the v8.1 air-pollutant release, not the 2025 GHG booklet the benchmark value comes from,
so the scope is \emph{inferred from EDGAR's category-resolved Pakistani source files}
rather than documented by the product being compared against. The distinction is not
pedantry---it is the same class of claim this paper elsewhere insists on labelling---and
it would be retired by EDGAR's per-category GHG product, which this study attempted to
fetch and could not obtain. On that inference EDGAR's Pakistani transport is road + rail
+ aviation, and the matched construction on this side is the same three components.

\begin{table*}[htbp]
\centering
\scriptsize
\setlength{\tabcolsep}{4pt}
\begin{tabular}{@{}p{0.17\linewidth}p{0.20\linewidth}p{0.18\linewidth}lp{0.24\linewidth}@{}}
\toprule
\textbf{Component} & \textbf{This paper} & \textbf{EDGAR ``transport''} &
\textbf{Period} & \textbf{Treatment in the comparison} \\
\midrule
Road (1A3b) & conditional residual, 40.19 Mt & 1.A.3.b present & FY2023--24 & compared \\
Railways (1A3c) & separate sector, 0.39 Mt & 1.A.3.c present & FY2023--24 & added \\
Domestic aviation (1A3a) & separate sector, 0.30 Mt & 1.A.3.a present & FY2023--24 &
  added, on the same fiscal year \\
Domestic navigation (1A3d) & OCAC Bunkers row, 3.1 kt $\approx$ 0.010 Mt (screened) &
  NO 1.A.3.d row for Pakistan & FY2023--24 & on neither side of the headline;
  all ten held editions converted in the navigation screen \\
Mixed off-road & 1.73 Mt, reported 1A4c & outside EDGAR transport (1A4/1A2) &
  FY2023--24 & excluded from headline; secondary line includes it \\
Bulk captive/industrial diesel & excluded by construction (Industry row, 400 kt) &
  in EDGAR Industrial Combustion & FY2023--24 & outside both transport totals \\
Retail agri-stationary + captive & inside road residual, unresolved & likely 1A4/1A2 &
  FY2023--24 & pushes this paper up; priced by the retail non-road transfers \\
Calendar convention & FY2023--24 (July--June) & calendar 2023 and 2024, blended &
  --- & primary comparison; single years and superseded rebuilds shown below \\
\bottomrule
\end{tabular}
\caption{The scope-and-period crosswalk, on EDGAR's own category evidence. Both sides
are AR5 \COe{}, all components on the FY2023--24 reference period against the blended
CY2023/CY2024 EDGAR columns. Because EDGAR's transport boundary is inferred from a
different release rather than published with the benchmark, every row here is a
scope-approximate diagnostic, not a validation.}
\label{tab:crosswalk}
\end{table*}

The headline: road residual (\SI{40.19}{\mega\tonne}) plus railways
(\SI{0.39}{\mega\tonne} strict 1A3c traction, FY2023--24---the non-traction
\SI{0.005}{\mega\tonne} is excluded because EDGAR's inferred scope has no such
category) plus domestic aviation (\SI{0.30}{\mega\tonne},
FY2023--24) gives \SI{40.88}{\mega\tonne}, or \num{0.85} of EDGAR's fiscal-year
blend of \SI{48.13}{\mega\tonne}. Against the 2024 column alone it is \num{0.84} and
against 2023 alone \num{0.86}---the width of that spread is what "fiscal year against
calendar year" costs, and it is smaller than the scope choices around it. Including the
mixed off-road class---which EDGAR books outside its transport aggregate---gives
\SI{42.62}{\mega\tonne}, \num{0.89}. Every construction is tabulated with its rank in
\texttt{outputs/national\_edgar\_transport\_comparison.csv}, including the superseded
rc4 construction above. The engine sits \emph{under} the benchmark in all of them, and
the retail-channel contamination that pushes this paper's residual the other way is
priced by the retail non-road transfers rather than asserted away. The gap between
\num{0.70} and \num{0.85} is not an engine change: it is what moving both sides of the
ratio onto the same period does, because the CY2021/CY2022 EDGAR window straddled the
demand peak that preceded the mid-2022 contraction.

One property of that comparison is worth naming because it is reassuring rather than
convenient: the off-road-inclusive construction is almost unmoved by the class-split
reweighting. Energy weighting moves diesel \emph{between} 1A3b and 1A4c, so a total
spanning both barely notices, which is what a partition should look like from outside.

One further construction is retained as history. On calendar-year 2022 controls,
rebuilt from the monthly panel (January--June direct editions plus the December
cumulative, which is July--December year-to-date by the panel's own tested identity), the
construction is about \num{0.82} of EDGAR's 2022 column: calendar 2022 straddles the
mid-2022 demand collapse, so that comparison is honest and wide at once
(\texttt{outputs/transport\_calendar\_year\_controls\_2024.csv}; provincial detail is
structurally unavailable for the calendar year's second half). It is ranked superseded
beside the rc4 blend, for the same reason. And the previous edition
sat \emph{above} EDGAR on a wrong-provenance control (Section~\ref{sec:data}); that
reversal, produced by reading the controls from their primary table, is itself evidence
that the earlier agreement was coincidence rather than validation.

\subsection{Internal closure checks}

Four properties are asserted on every run and can fail. Closure is an arithmetic property
of the ledger, not evidence that the residual is road; the evidence for the residual's
composition is the customer-class semantics of the control row, and its remaining
contamination is stated in Table~\ref{tab:ledger} rather than hidden by the arithmetic.

\textbf{The shares sum to one.} $\sum_c s_{c,f} = 1$ for each fuel, so the class emissions sum
to the control total times the factor. This is what makes the anchoring claim true rather than
approximately true.

\textbf{The off-road memo is empty.} The previous treatment carried off-road as a \COtwo{}-only
memo item outside all totals. Now that it is a reported category, a test asserts the memo
contains nothing, so the same diesel cannot appear as both.

\textbf{All four ledgers close.} Opening equals claims plus closing residual for each
(fuel, tier) ledger, to $10^{-9}$ relative; the run aborts otherwise. The unresolved
claimants carry no quantity by construction, so nothing can be invented into the closure.

\textbf{The rail figures agree without being equal.} The year-matched Year Book row and
the OCAC Railways row must sit within tolerance of each other and must \emph{not}
coincide. The inequality is a regression guard against the two independently sourced
fields accidentally collapsing to one stored value; their independence is established by
their separate reporting chains, not by the test.

\subsection{The fleet-implied fuel demand, against the observed controls}
\label{sec:fuelgap}

The comparison the anchoring argument rests on can now be shown rather than asserted. A
conventional bottom-up inventory would compute each fuel's demand from the fleet:
$F_f = \sum_c \mathrm{vkt}_{c,f} \times e_{c,f}$. Using this paper's own
vehicle-kilometres and \emph{declared engineering bands} for the per-class energy
intensities $e_{c,f}$ (judgement values, not Pakistani measurements---the same epistemic
class as the CNG check's plausibility band,
\texttt{outputs/transport\_\allowbreak cng\_\allowbreak crosscheck\_\allowbreak 2024.csv},
against the metered series in
\texttt{outputs/transport\_\allowbreak cng\_\allowbreak 2024.csv}), the bottom-up build
misses the observed controls in both directions:

\begin{table}[htbp]
\centering
\footnotesize
\begin{tabular}{@{}lrrrr@{}}
\toprule
\textbf{Fuel} & \textbf{Fleet-implied (TJ)} & \textbf{Observed (TJ)} &
\textbf{Ratio} & \textbf{Verdict} \\
\midrule
Petrol & 208,714--374,877 & 316,882 & 0.84 & fleet-model low \\
Diesel & 112,392--205,316 & 238,712 & 0.59 & fleet-model low \\
CNG    & 54,459--98,535   & 21,488  & 3.05 & fleet-model high \\
\bottomrule
\end{tabular}
\caption{Fleet-implied fuel demand at the declared intensity bands (low--high; the ratio
is band-midpoint over control) against the observed controls
(\texttt{outputs/transport\_fleet\_fuel\_diagnostic\_2024.csv}). Even the high band
undershoots the diesel control, while even the low band overshoots metered CNG: a
bottom-up build on this fleet data would err by a factor of two in one direction and
three in the other. The diagnostic flags and never adjusts.}
\label{tab:fuelgap}
\end{table}

This is the demonstration behind the paper's central design choice. Fuel observed at the
point of sale sets the magnitude; the fleet model, which cannot even reproduce the fuel,
is confined to distributing it.

One caveat on this diagnostic is owed to the reader, because this edition created it.
The intensity bands used here are now also the weights of the class split
(Section~\ref{sec:methods}), so the diagnostic is no longer independent of the
composition it is quoted beside---both read one parameter set, deliberately, so the two
cannot silently diverge. What the diagnostic still tests independently is the
\emph{magnitude}: the class split is normalised and therefore cannot reproduce a fuel
total at all, whatever intensities it is given. That is the comparison in the table, and
the reweighting did not change it by a tonne.

\subsection{What is not validated}

The class split has no independent check at all. No nationally representative Pakistani
vehicle-kilometre survey has been identified for this study, and no Pakistani on-road
fuel-economy measurement either, so the composition in Section~\ref{sec:results} rests on
registration counts, assumed annual kilometres and declared energy intensities---three
inputs, none of them measured in the country. This is stated here rather than in the
limitations alone because a reader may reasonably want the class table and should know
that it is the least supported thing in the paper. It is also why the class results are
reported as a bracketed range: the composition scenarios of
Section~\ref{sec:results} move the two-wheeler share of on-road \COtwo{} across roughly a
quarter of its own value without moving the sector total at all.

\subsection{CNG on two activity bases}
\label{sec:cngcheck}

CNG is the one fuel in this sector for which two independent descriptions of the same vehicles
exist. The emissions come from OGRA's metered gas. The geography and the air-pollutant layer
come from the fleet model, which reports CNG vehicle-kilometres by class. Neither is derived
from the other, so dividing one by the other tests whether they describe the same fleet---the
same move the power sector makes with its implied heat rate.

The metered gas is \SI{21488}{} net terajoules---the reported volume converted at a gross
heating value of \SI{1000}{} Btu/scf and then to a net basis at \num{0.90}, the IPCC
convention, with both parameters pinned equal to the captive-power category's by test. The
fleet model puts \SI{8.55}{\percent} of road vehicle-kilometres on CNG, \SI{2.36e10}{} of
them. Together, implying \SI{0.912}{\mega\joule\per\kilo\metre}.

That is low. Published consumption rates put a CNG three-wheeler nearer
\SI{1.6}{\mega\joule\per\kilo\metre} and a CNG car nearer \SI{3}{}, and
\SI{74}{\percent} of this modelled travel is three-wheelers, so the metered gas is well short
of what the modelled distance would require. The two bases disagree, and the direction says the
fleet model overstates CNG travel.

Three qualifications. The per-vehicle consumption rates behind that comparison are generic
engineering figures rather than Pakistani measurements, so the check flags a disagreement and
never adjusts anything. The disagreement is confined to what rides vehicle-kilometres---the
provincial split and the air-pollutant layer---because the emissions come from the gas. And it
is not a new weakness: this sector already records that its class split rides vehicle-kilometres
and over-weights small vehicles, and already argues that magnitude must therefore come from
metered fuel. CNG is simply the first category where an independent fuel quantity exists to test
that argument, and the test agrees with it.
